\documentclass[11pt]{article}
\usepackage{rabbithole}

\phasenumber{5}
\phasetitle{Rewriting history}
\phasesubtitle{Interactive rebase, and the difference between what happened and what you want to publish.}

\hypersetup{pdftitle={Rabbit Holes, Phase 5: Rewriting history},
            pdfauthor={Accelerate, Manipal University Jaipur},
            pdfsubject={git and GitHub handbook}}

\begin{document}
\maketitlepage

\section{Two different histories}

The history of how you actually worked is a mess. You committed something called
\co{wip}, then \co{wip 2}, then \co{fix typo}, then \co{actually fix typo}, then you
realised the whole approach was wrong and started again.

The history you want a reviewer to read is a short sequence of complete, sensible
changes, each of which does one thing and has a message explaining why.

Those two things are almost never the same, and git is one of very few version
control systems that lets you turn the first into the second before anyone else
sees it. That is what this phase is about.

\begin{why}
This is not vanity. A clean history is a working tool. \co{git bisect} can find the
commit that broke your build, but only if each commit builds. \co{git revert} can
undo a feature cleanly, but only if the feature is not smeared across nine commits
mixed with unrelated work. \co{git blame} tells you why a line exists, but only if
the commit it points at has a real message.

Phase 6 is entirely about tools that only work well on a tidy history. This phase
is how the history gets tidy.
\end{why}

\section{The golden rule}

\begin{trap}
\textbf{Do not rewrite history that other people have.}

Rewriting produces new commits with new hashes. Anyone whose branch is built on the
old commits now has work attached to objects that no longer exist upstream, and
sorting that out is their problem, caused by you.

In practice: rewrite freely on your own unpushed branch. Rewrite carefully on a
pushed branch that only you are working on, which is what \co{--force-with-lease}
is for. Never rewrite \co{main}, or any branch other people build on, unless the
whole team has agreed and is standing by.
\end{trap}

Everything below is safe on a branch that is yours alone. That covers the large
majority of real use.

\section{Interactive rebase}

The main tool.

\begin{shell}
git rebase -i HEAD~5
\end{shell}

This says: take the last five commits, and let me edit the plan for replaying them.
Your editor opens with a to-do list:

\begin{terminal}
pick 3a1f2b9 Add poem parser
pick 7c4d8e1 wip
pick 2b9a0c3 wip 2
pick 8e1f4a7 fix typo
pick 5d3c2b8 Handle empty input

# Rebase 9f2e1d0..5d3c2b8 onto 9f2e1d0 (5 commands)
\end{terminal}

Two things surprise everyone the first time:

\begin{enumerate}
  \item \textbf{The list is oldest first.} \co{git log} shows newest first. The
    rebase to-do shows the order commits will be replayed in, which is the opposite.
  \item \textbf{You edit the list, not the commits.} You are writing a script. When
    you save and close, git runs it.
\end{enumerate}

Change the word at the start of a line to change what happens to that commit:

\vspace{0.5em}
\begin{tabularx}{\linewidth}{@{}L{3.2cm} X@{}}
\toprule
\rhtablehead{Command} & \rhtablehead{Effect} \\
\midrule
\co{pick} & Keep the commit as it is. The default. \\
\co{reword} & Keep the changes, edit the message. \\
\co{edit} & Stop here so you can change the commit's contents. \\
\co{squash} & Merge into the previous commit, and combine both messages. \\
\co{fixup} & Merge into the previous commit, and throw this message away. \\
\co{drop} & Delete the commit entirely. Deleting the line does the same. \\
\co{exec} & Run a shell command at this point. \\
\co{break} & Stop here for no reason other than that you asked. \\
\bottomrule
\end{tabularx}
\vspace{0.5em}

You can also \textbf{reorder} commits by moving lines around, which git will do
without complaint as long as the changes do not depend on each other.

\subsection{Squashing a messy branch}

The most common use by a wide margin. Turn five commits into one:

\begin{terminal}
pick   3a1f2b9 Add poem parser
fixup  7c4d8e1 wip
fixup  2b9a0c3 wip 2
fixup  8e1f4a7 fix typo
squash 5d3c2b8 Handle empty input
\end{terminal}

The three \co{wip} messages are discarded by \co{fixup}. The last one is
\co{squash}, so git will open an editor letting you write a combined message. The
result is one commit called something like ``Add poem parser, handling empty
input''.

\begin{note}
\co{squash} and \co{fixup} both merge upward, into the line above. That is why the
first line must stay \co{pick}: there is nothing above it to merge into.
\end{note}

\subsection{Rewording without touching content}

\begin{terminal}
pick   3a1f2b9 Add poem parser
reword 7c4d8e1 asdfgh
pick   2b9a0c3 Handle empty input
\end{terminal}

Git replays the first commit, stops at the second, opens your editor with the old
message, and continues once you save.

\subsection{Splitting one commit into two}

You committed two unrelated changes together and want them separated. Mark it
\co{edit}:

\begin{terminal}
pick 3a1f2b9 Add poem parser
edit 7c4d8e1 Add parser and fix unrelated bug
\end{terminal}

Git stops with that commit checked out. Undo it, keeping the changes on disk:

\begin{shell}
git reset HEAD~1
\end{shell}

Now stage and commit the two halves separately:

\begin{shell}
git add parser.py
git commit -m "Add poem parser"
git add utils.py
git commit -m "Fix off-by-one in line counter"
git rebase --continue
\end{shell}

\begin{tryit}
If the two changes are in the \emph{same file}, stage them separately with:

\begin{shell}
git add -p utils.py
\end{shell}

Git walks you through each chunk of the diff and asks whether to stage it.
Press \co{y} to stage, \co{n} to skip, \co{s} to split the chunk into smaller ones,
and \co{?} for the full list.

\co{git add -p} is one of the highest value commands in git and almost nobody is
taught it. It is what makes small, focused commits practical instead of aspirational.
\end{tryit}

\section{The autosquash workflow}

Interactive rebase is powerful and it makes you stop and edit a to-do list. There is
a faster path for the specific case of ``I need to fix a commit I made earlier on
this branch''.

Find the commit you need to amend, then:

\begin{shell}
git add the-fix.py
git commit --fixup 3a1f2b9
\end{shell}

This makes an ordinary commit with the message \co{fixup! Add poem parser}. Keep
working, make as many of these as you like. When you are ready to tidy up:

\begin{shell}
git rebase -i --autosquash HEAD~10
\end{shell}

Git has already rearranged the to-do list for you: each \co{fixup!} commit has been
moved directly beneath the commit it refers to and marked \co{fixup}. Usually you
save the file without changing anything.

Make it the default so you never have to remember the flag:

\begin{shell}
git config --global rebase.autoSquash true
\end{shell}

\begin{why}
This exists because the alternative is worse. Without it you either interrupt your
work to do a full interactive rebase every time you notice a small mistake, or you
push a branch with six commits called ``address review comments'', which tells a
future reader nothing at all.
\end{why}

\section{Cherry-pick}

Take one commit from anywhere and apply it here.

\begin{shell}
git cherry-pick 8e1f4a7
\end{shell}

Useful when a bug fix landed on a feature branch and needs to go to \co{main} now,
without the rest of the feature.

\begin{shell}
git cherry-pick 8e1f4a7 2b9a0c3
git cherry-pick main~5..main~2
git cherry-pick -n 8e1f4a7
\end{shell}

\co{-n} applies the change without committing, so you can adjust it first.

\begin{trap}
Cherry-picking makes a \emph{copy}, with a different hash. If you cherry-pick from a
branch and later merge that same branch, git may present the same change twice and
conflict with itself.

For one urgent fix, cherry-pick. For anything ongoing, merge properly instead.
Repeated cherry-picking between long-lived branches produces duplicated history that
becomes genuinely hard to reason about.
\end{trap}

\section{rebase -{}-onto}

The surgical version, and the one worth knowing when you have made a branch from the
wrong place.

You branched \co{feature} off \co{experimental}, then realised it should have come
off \co{main}. You want your three commits, without dragging \co{experimental}
along.

\begin{shell}
git rebase --onto main experimental feature
\end{shell}

Read it as three parts: \textbf{onto} where they should land, \textbf{from} the old
base to exclude, \textbf{which} branch to move.

\begin{terminal}
before                         after

A---B---C  main                A---B---C  main
     \                                  \
      D---E  experimental                F'--G'--H'  feature
           \
            F---G---H  feature
\end{terminal}

Commits \co{D} and \co{E} are left behind. \co{F}, \co{G} and \co{H} are replayed
on top of \co{main} as new commits.

\section{Rewriting all of history}

Sometimes the problem is not the last few commits. Somebody committed a password
three months ago, or a 200 MB video file that has been slowing every clone since.

Removing the file in a new commit is not enough. It is still in the history, still
in every clone, and a secret in git history should be treated as public.

The tool is \co{git-filter-repo}, which is separate from git and needs installing:

\begin{shell}
pip install git-filter-repo
\end{shell}

Remove a file from every commit that ever contained it:

\begin{shell}
git filter-repo --path secrets.env --invert-paths
\end{shell}

Or purge a string everywhere it appears:

\begin{shell}
git filter-repo --replace-text expressions.txt
\end{shell}

\begin{trap}
This rewrites every commit hash in the repository. Every open pull request breaks,
every clone becomes incompatible, and everyone must re-clone from scratch. It is a
coordinated operation, not something to do on a Tuesday afternoon.

And it does not un-leak the secret. Anyone who cloned before you rewrote still has
it, and GitHub may have cached it. \textbf{Rotate the credential first, then clean
the history.} In that order, always. The rewrite is housekeeping, the rotation is
the actual fix.
\end{trap}

\begin{note}
You will see \co{git filter-branch} in older documentation. It is slow, easy to get
wrong, and git itself now prints a warning recommending you use
\co{git-filter-repo} instead. The BFG Repo-Cleaner is another good option, faster
than \co{filter-branch} and simpler for the common cases.
\end{note}

\section{When a rebase goes wrong}

It will, and it is fine.

\vspace{0.5em}
\begin{tabularx}{\linewidth}{@{}L{5.8cm} X@{}}
\toprule
\rhtablehead{Command} & \rhtablehead{What it does} \\
\midrule
\co{git rebase -{}-abort} & Stop, put everything back the way it was \\
\co{git rebase -{}-skip} & Discard the commit being applied, carry on \\
\co{git rebase -{}-continue} & Carry on after resolving a conflict \\
\co{git rebase -{}-edit-todo} & Change the plan mid-rebase \\
\co{git status} & Tells you where in the rebase you are \\
\bottomrule
\end{tabularx}
\vspace{0.5em}

\co{--abort} works at any point and always returns you to exactly where you started.
There is no state a rebase can leave you in that you cannot back out of.

If you have already finished a bad rebase:

\begin{shell}
git reset --hard ORIG_HEAD
\end{shell}

Or find the pre-rebase position in \co{git reflog} and reset to it. Phase 4 covers
this in full.

\begin{deeper}
Since Git 2.26 the default rebase engine is the ``merge'' backend rather than the
old patch based one. This is why rebase now handles renames, empty commits and
directory moves far better than the reputation it earned a decade ago.

It also means \co{git rebase} and \co{git merge} share conflict resolution
machinery, so \co{rerere} works with both. Phase 7 covers \co{rerere}, which
remembers how you resolved a conflict and replays that resolution automatically the
next time it sees the same one. On a long-lived branch that you rebase repeatedly,
it is transformative.
\end{deeper}

\section{Reference}

\vspace{0.5em}
\begin{tabularx}{\linewidth}{@{}L{6.4cm} X@{}}
\toprule
\rhtablehead{Goal} & \rhtablehead{Command} \\
\midrule
Edit the last commit & \co{git commit -{}-amend} \\
Tidy the last N commits & \co{git rebase -i HEAD\textasciitilde N} \\
Squash a branch into one commit & \co{rebase -i}, then \co{fixup} every line but the first \\
Fix an earlier commit, later & \co{git commit -{}-fixup <sha>} then \co{rebase -i -{}-autosquash} \\
Split a commit in two & \co{edit}, then \co{git reset HEAD\textasciitilde 1}, then commit twice \\
Stage part of a file & \co{git add -p} \\
Copy one commit here & \co{git cherry-pick <sha>} \\
Move a branch to a new base & \co{git rebase -{}-onto <new> <old> <branch>} \\
Remove a file from all history & \co{git filter-repo -{}-path <f> -{}-invert-paths} \\
Undo a finished rebase & \co{git reset -{}-hard ORIG\_HEAD} \\
\bottomrule
\end{tabularx}
\vspace{0.5em}

\checkyourself{
\begin{enumerate}
  \item Why is the interactive rebase to-do list in the opposite order to
    \co{git log}?
  \item What is the difference between \co{squash} and \co{fixup}?
  \item Walk through splitting one commit into two.
  \item What does \co{git commit -{}-fixup} buy you over doing a rebase immediately?
  \item You leaked an API key in a commit from March. What are the two steps, and
    which comes first?
\end{enumerate}}

\vspace{1em}
{\color{rhborder}\rule{\linewidth}{0.8pt}}

\textbf{Next:} Phase 6, \emph{Forensics}, which is about reading history rather than
writing it: finding the commit that broke the build automatically, searching every
version of every file for a string that no longer exists, and working out who to ask.

\end{document}
